\documentclass[12pt]{article}

\usepackage[utf8]{inputenc}
\usepackage[T1]{fontenc}
\usepackage[spanish]{babel}
\usepackage[margin=2.5cm]{geometry}
\usepackage{listings}

\title{Preparación del entorno de desarrollo en macOS}
\author{Daniel Soto Villada}
\date{\today}

\lstset{
  basicstyle=\ttfamily\small,
  breaklines=true,
  frame=single,
  keepspaces=true,
  columns=fullflexible
}

\begin{document}

\maketitle

Esta guía reúne los pasos iniciales para preparar un entorno de desarrollo
en macOS: instalar las herramientas de línea de comandos de Xcode, comprobar
la disponibilidad de Git e instalar Homebrew.

\section{Instalar las herramientas de línea de comandos de Xcode}

Abre Terminal y ejecuta:

\begin{lstlisting}
xcode-select --install
\end{lstlisting}

Sigue las instrucciones que aparecen en pantalla para completar la instalación.

\section{Comprobar Git}

Comprueba si Git está disponible con el siguiente comando:

\begin{lstlisting}
git --version
\end{lstlisting}

Si Git no está instalado, instala las herramientas de línea de comandos de
Xcode siguiendo el paso anterior y vuelve a ejecutar el comando.

\section{Instalar Homebrew}

Cuando las herramientas de línea de comandos estén instaladas, ejecuta en
Terminal el instalador de Homebrew:

\begin{lstlisting}
bash -c "$(curl -fsSL https://raw.githubusercontent.com/Homebrew/install/HEAD/install.sh)"
\end{lstlisting}

Sigue las indicaciones del instalador para finalizar la configuración.

\section{¿Qué es Homebrew y para qué sirve?}

Homebrew es un gestor de paquetes para macOS que permite instalar y actualizar
programas desde Terminal con comandos sencillos. No es imprescindible para
programar, pero facilita la preparación del entorno y la instalación de
herramientas y componentes para programación y análisis de datos, como Python,
R y Jupyter.

\end{document}
